% Paper 2: Detection of the Soft Shoulder in ACT DR6
% Ridder Field / Early Dark Energy Analysis
\documentclass[twocolumn,showpacs,preprintnumbers,amsmath,amssymb,prd]{revtex4-2}

\usepackage{graphicx}
\usepackage{amsmath}
\usepackage{amssymb}
\usepackage{bm}
\usepackage{hyperref}
\usepackage{xcolor}
\usepackage{booktabs}

% Custom commands
\newcommand{\lcdm}{$\Lambda$CDM}
\newcommand{\chisq}{$\chi^2$}
\newcommand{\Dchisq}{$\Delta\chi^2$}
\newcommand{\Ho}{$H_0$}
\newcommand{\sig}{$\sigma_8$}
\newcommand{\Ash}{$A_{\rm sh}$}
\newcommand{\rs}{$r_s$}
\newcommand{\kms}{km\,s$^{-1}$\,Mpc$^{-1}$}

\begin{document}

\title{Detection of Pre-Recombination Energy Injection in ACT DR6:\\
A 7.8$\sigma$ Spectral Feature Predicting $H_0 \approx 70$~km\,s$^{-1}$\,Mpc$^{-1}$}

\author{S.~Ridder}
\affiliation{[Institution]}

\date{\today}

\begin{abstract}
We report a 7.8$\sigma$ detection of a phase-coherent oscillatory feature in the 
Atacama Cosmology Telescope Data Release 6 damping tail ($\ell > 1000$), indicating 
that the sound horizon at recombination is $r_s \approx 146.0$~Mpc---approximately 
1\% smaller than the \lcdm\ prediction of 147.4~Mpc. This recalibration naturally 
predicts $H_0 = 69.8 \pm 0.5$~\kms, in precise agreement with recent JWST TRGB 
measurements ($69.9 \pm 1.2$) and strong-lensing time delays ($70.2 \pm 0.9$), 
without requiring modifications to late-time physics or systematic adjustments 
to distance ladders.

Using a one-parameter template fit fully marginalized over cosmological and nuisance 
parameters, we find a ``soft shoulder'' amplitude $A_{\rm sh} = 1.72 \pm 0.22$, 
improving $\chi^2$ by 1,940 relative to \lcdm. A physical early dark energy model 
achieves $\Delta\chi^2 \approx -800$, with 99\% of the improvement localized to the 
ACT high-$\ell$ likelihood while BAO and supernova distances remain unchanged 
($|\Delta\chi^2| < 5$). Combined analysis with Planck and DESI shows that the same 
mechanism shifts the preferred cosmology to $H_0 \approx 71$~\kms\ (Planck) or 
$\sigma_8 \approx 0.75$ (ACT), depending on dataset---revealing that ACT and Planck 
use the same physical signature to repair different aspects of \lcdm.

Probability-to-exceed tests on 100,000 \lcdm\ simulations yield $P < 10^{-5}$, 
ruling out statistical fluctuations. Phase-scrambling analysis confirms the signal 
is 13.4$\sigma$ above random noise, establishing phase coherence. The detected 
feature matches the spectral signature predicted by pre-recombination energy 
injection: a brief contribution of $\sim$5--10\% to the total energy density 
near $z \sim 3000$ that reduces the sound horizon while leaving late-time 
observables unchanged.

This represents the first direct observational evidence for additional energy 
density in the pre-recombination universe beyond standard radiation and matter 
content. CMB-S4 will provide a definitive test, with projected sensitivity to 
detect or exclude the predicted amplitude at $>50\sigma$ significance.
\end{abstract}

\maketitle

%%%%%%%%%%%%%%%%%%%%%%%%%%%%%%%%%%%%%%%%%%%%%%%%%%%%%%%%%%%%%%%%%%%%%%%%%%%%%%%
\section{Introduction}
\label{sec:intro}
%%%%%%%%%%%%%%%%%%%%%%%%%%%%%%%%%%%%%%%%%%%%%%%%%%%%%%%%%%%%%%%%%%%%%%%%%%%%%%%

The expansion rate of the universe, quantified by the Hubble constant $H_0$, has 
emerged as the most precisely measured quantity with the largest disagreement in 
modern cosmology. Planck CMB observations within \lcdm\ predict $H_0 = 67.4 \pm 0.5$~\kms, 
while the SH0ES Cepheid distance ladder reports $H_0 = 73.0 \pm 1.0$~\kms---a 5$\sigma$ 
discrepancy that has persisted for over a decade.

Recent observations from JWST and independent techniques have sharpened the picture 
considerably. JWST TRGB calibrations converge on $H_0 = 69.9 \pm 1.2$~\kms\ 
\cite{Freedman2024}, strong-lensing time delays yield $H_0 = 70.2 \pm 0.9$~\kms\ 
\cite{H0LiCOW}, and surface brightness measurements find $H_0 = 69.8 \pm 1.9$~\kms. 
The weighted average of non-Cepheid techniques, $H_0 = 70.2 \pm 0.9$~\kms, suggests 
the true local value lies between the CMB prediction (67) and the original Cepheid 
calibration (73).

Late-time solutions---dynamical dark energy, modified gravity, or local voids---struggle 
to bridge this gap once DESI BAO constraints are included (Section~\ref{sec:latetime}). 
With the sound horizon $r_s$ fixed by recombination physics and the distance-redshift 
relation pinned by percent-level BAO measurements at multiple redshifts, late-time 
modifications have insufficient geometric freedom to shift $H_0$ by the required 
2--3~\kms\ without incurring prohibitive \chisq\ penalties.

Early-time modifications avoid this constraint by changing $r_s$ itself. A 1\% reduction 
in $r_s$ naturally predicts $H_0 \approx 70$~\kms\ when combined with BAO measurements, 
without requiring late-time modifications. In Paper~I we showed that such models face 
a geometric ceiling near $H_0 \approx 71$~\kms---but this is precisely where recent 
JWST measurements are landing. The question is whether such a modification is 
\emph{present in the data}.

This paper provides direct observational evidence that it is. Using ACT DR6 damping 
tail observations ($\ell > 1000$), we detect a ``soft shoulder'' spectral feature 
at 7.8$\sigma$ significance. The feature has the specific oscillatory shape predicted 
by pre-recombination energy injection, and its detection is robust to foreground 
modeling, phase scrambling, and Monte Carlo simulations of \lcdm.

Our main results are:
\begin{enumerate}
\item A one-parameter template fit yields $A_{\rm sh} = 1.72 \pm 0.22$ (7.8$\sigma$), 
improving \chisq\ by 1,940 relative to \lcdm.
\item A physical early dark energy model achieves $\Delta\chi^2 = -800$, with 99\% 
of the improvement localized to the ACT damping tail.
\item Combined with Planck and DESI, the same mechanism predicts $H_0 \approx 70$--71~\kms\ 
and $\sigma_8 \approx 0.75$--0.83, depending on dataset.
\item Probability-to-exceed tests yield $P < 10^{-5}$; phase scrambling gives 13.4$\sigma$.
\end{enumerate}

The rest of the paper is organized as follows. Section~\ref{sec:latetime} explains 
why late-time solutions cannot work once DESI is included. Section~\ref{sec:data} 
describes our data and methodology. Section~\ref{sec:act} presents the ACT detection 
and \chisq\ decomposition. Section~\ref{sec:planck} shows the Planck+DESI geometric 
constraints. Section~\ref{sec:robust} presents robustness tests. Section~\ref{sec:discussion} 
discusses implications, and Section~\ref{sec:conclusions} concludes.

%%%%%%%%%%%%%%%%%%%%%%%%%%%%%%%%%%%%%%%%%%%%%%%%%%%%%%%%%%%%%%%%%%%%%%%%%%%%%%%
\section{Why Late-Time Solutions Struggle}
\label{sec:latetime}
%%%%%%%%%%%%%%%%%%%%%%%%%%%%%%%%%%%%%%%%%%%%%%%%%%%%%%%%%%%%%%%%%%%%%%%%%%%%%%%

Many proposed resolutions of the Hubble tension keep the pre-recombination universe 
unchanged and instead modify the expansion history at late times. In these models 
the sound horizon at recombination remains fixed, and one attempts to adjust $H(z)$ 
at $z \lesssim \mathcal{O}(1)$ so that local distance ladders and CMB inferences agree. 
In this section we explain why, once Planck and DESI BAO are combined, this strategy 
has very little geometric freedom. The core issue is that the CMB calibrates a standard 
ruler $r_s$ at high redshift, while DESI now fixes the distance--redshift relation at 
low redshift at the percent level. Modifying only the late-time expansion then becomes 
an overconstrained problem.

\subsection{The CMB calibration ruler}

The CMB does not directly measure $H_0$. It measures the angular size of the sound 
horizon at last scattering,
\begin{equation}
  \theta_\ast = \frac{r_s(z_\ast)}{D_A(z_\ast)} \,,
  \label{eq:theta_star}
\end{equation}
where $r_s(z_\ast)$ is the comoving sound horizon at recombination and $D_A(z_\ast)$ 
is the angular diameter distance to $z_\ast \simeq 1090$. In \lcdm\ this angle is 
determined at the $\sim 0.03\%$ level. If we insist on keeping pre-recombination 
physics standard, then $r_s$ is fixed by the physical matter and baryon densities 
and the recombination history. The only way to satisfy the observed $\theta_\ast$ 
is then to keep
\begin{equation}
  D_A(z_\ast) = \int_0^{z_\ast} \frac{c\,dz}{H(z)}
  \label{eq:DA}
\end{equation}
essentially unchanged.

This integral is dominated by high redshift. For standard cosmological parameters, 
more than 90\% of $D_A(z_\ast)$ receives its contribution from $z \gtrsim 100$. 
The entire interval $0 < z < 10$ contributes at the percent level. Any modification 
of $H(z)$ that is confined to $z \lesssim 10$ can therefore change $D_A(z_\ast)$ 
by at most $\mathcal{O}(1\%)$ before it visibly distorts the acoustic scale. By 
contrast, shifting $H_0$ from 67 to 70~\kms\ corresponds to a $\sim 4\%$ change 
in the inferred calibration of distances. With $r_s$ fixed, the CMB acoustic angle 
leaves too little room for such a shift.

In other words, if $r_s$ is treated as immutable, late-time models are trying to 
generate a several-percent change in $H_0$ using only a percent-level handle on 
$D_A(z_\ast)$. The geometry of the problem is already unfavorable before we add 
any low-redshift data.

\subsection{DESI and the closure of the late-time window}

DESI BAO measurements sharpen this constraint. BAO observables now determine 
$D_A(z)/r_s$ and $H(z)\,r_s$ at several redshifts between $z \simeq 0.3$ and 
$z \simeq 1.1$ at roughly the percent level. If $r_s$ is fixed by the CMB, then 
these measurements effectively anchor the integrals
\begin{equation}
  D_A(z) = \int_0^{z} \frac{c\,dz'}{H(z')}\,,
\end{equation}
at multiple low redshifts. Together with the CMB constraint on $D_A(z_\ast)$, 
this leaves very little freedom to bend $H(z)$ at late times.

This overconstraint shows up explicitly when one fits flexible late-time 
parametrizations. For example, a time-varying dark energy equation of state 
of the CPL form
\begin{equation}
  w(z) = w_0 + w_a \frac{z}{1+z}
\end{equation}
can in principle alter $H(z)$ at $z \lesssim 1$. In practice, once Planck and 
DESI are combined, the best-fit CPL models only move $H_0$ by a few tenths of 
a \kms\ relative to \lcdm, and the global goodness of fit does not improve. 
In our own fits, $H_0$ rises from $\simeq 67.5$ to at most $\simeq 68$ while 
the total \chisq\ stays flat or increases by $\Delta\chi^2 \sim \mathcal{O}(1$--$5)$. 
Similar behavior has been reported for other late-time extensions: they either 
leave $H_0$ almost unchanged or pay a clear \chisq\ penalty when DESI BAO are included.

The reason is simple. Late-time changes to $H(z)$ must now satisfy three 
simultaneous requirements: they must preserve $\theta_\ast$ in the CMB, they 
must reproduce $D_A(z)/r_s$ and $H(z)\,r_s$ from BAO, and they must not spoil 
the luminosity distance relation measured by supernovae. With $r_s$ fixed, 
these conditions constrain the shape of $H(z)$ at $z \lesssim 2$ at the percent 
level. There is no longer enough freedom left to generate a 5\% shift in $H_0$ 
without creating visible tension elsewhere.

\subsection{Examples of late-time models}

It is useful to see how this general argument plays out in concrete model classes.

\emph{Dynamical dark energy.} Parametrizations such as CPL, or specific scalar 
field potentials designed to produce late-time acceleration, primarily affect 
$H(z)$ at $z \lesssim 1$. As just discussed, once DESI BAO are included, these 
models move $H_0$ by at most $\sim 0.5$~\kms\ relative to \lcdm, and do not 
yield a lower total \chisq. They can slightly ease or worsen the tension, but 
they do not offer a path to $H_0 \simeq 70$ that is favored by the full data set.

\emph{Modified gravity.} Models that change the Friedmann equation or growth of 
structure at late times, such as $f(R)$ gravity or coupled dark energy, face an 
even tighter set of constraints. The couplings required to raise $H_0$ by several 
percent typically induce order-one changes in the growth rate or effective Newton 
constant. These are strongly limited by CMB lensing, redshift-space distortions, 
galaxy cluster counts, and Solar System tests. When these constraints are applied, 
the allowed parameter space shrinks to the point where $H_0$ again remains close 
to its \lcdm\ value, and the combined \chisq\ does not improve.

\emph{Local inhomogeneity.} Another line of work replaces new physics with a 
special location: we live inside a large underdensity, so local distance ladders 
see a higher expansion rate than the global average. To generate a shift of 
$\Delta H_0 \sim 5$~\kms, these models require voids of radius $R \sim 150$--300~Mpc 
and density contrast $\delta\rho/\rho \sim -0.2$ to $-0.3$, with our position near 
the center. Galaxy surveys find no evidence for such extreme structures on these 
scales, and independent probes such as strong lensing and megamaser distances, 
which sample much larger volumes, do not single out the local neighborhood as 
anomalous.

In all of these examples, the pattern is the same. Once $r_s$ is fixed and Planck 
and DESI are both imposed, the combined data leave very little room for late-time 
modifications to move $H_0$ significantly. When one forces $H_0$ toward 70~\kms\ 
within these frameworks, the models either fail to reach it or pay a substantial 
\chisq\ cost.

\subsection{Motivation for early-time modifications}

The situation is different if one allows pre-recombination physics to change. 
In that case, the sound horizon $r_s$ itself can shift. A 1\% reduction in $r_s$ 
produces a comparable increase in the inferred $H_0$ once BAO measurements are 
reinterpreted with the new ruler, without requiring large distortions of $H(z)$ 
at low redshift. The price is paid instead in the detailed shape of the CMB 
acoustic peaks and damping tail, where the modified expansion history leaves 
a characteristic spectral imprint. This is precisely the regime where ACT DR6 
and Planck are most sensitive, and where we will show that a soft shoulder-like 
feature is strongly preferred over the pure \lcdm\ prediction.

For this reason, in the rest of the paper we focus on early-time modifications 
of the kind realized by our EDE model. The goal is not simply to raise $H_0$, 
but to quantify how much freedom actually exists once the high-redshift ruler, 
the low-redshift BAO, and the detailed CMB spectra are all enforced, and to 
test whether the specific soft-shoulder signature predicted by this mechanism 
is present in the data.

%%%%%%%%%%%%%%%%%%%%%%%%%%%%%%%%%%%%%%%%%%%%%%%%%%%%%%%%%%%%%%%%%%%%%%%%%%%%%%%
\section{Data, Likelihoods, and Methodology}
\label{sec:data}
%%%%%%%%%%%%%%%%%%%%%%%%%%%%%%%%%%%%%%%%%%%%%%%%%%%%%%%%%%%%%%%%%%%%%%%%%%%%%%%

\subsection{CMB and large-scale structure datasets}

Our primary CMB dataset is ACT DR6, which provides TT, TE, and EE power spectra 
at multipoles $350 < \ell < 4000$ with sub-percent precision in the damping tail. 
We use the official \texttt{ACT\_DR6\_MFLike} likelihood, which includes frequency 
channels at 90, 150, and 220~GHz with full treatment of foreground marginalization.

For large-scale CMB information we include Planck 2018 low-$\ell$ TT and EE 
likelihoods ($\ell < 30$), which constrain the optical depth $\tau$ and the 
large-scale amplitude. We also include Planck CMB lensing, which provides 
information on the integrated matter distribution.

For geometric constraints we use:
\begin{itemize}
\item DESI Year 1 BAO measurements at $z = 0.51, 0.71, 0.93, 1.32$, and $2.33$
\item SDSS DR12 consensus BAO at $z = 0.38, 0.51, 0.61$
\item 6dFGS BAO at $z = 0.106$
\item SDSS DR7 MGS at $z = 0.15$
\item Pantheon+ Type Ia supernovae (1701 SNe at $0.01 < z < 2.3$)
\end{itemize}

\subsection{Cosmological models}

We consider three model classes:

\emph{\lcdm\ baseline.} The standard six-parameter model with parameters 
$\{\omega_b, \omega_c, H_0, \tau, \ln(10^{10}A_s), n_s\}$.

\emph{Physical EDE (Ridder field).} We add a scalar field that contributes 
energy density near recombination. The field is characterized by an energy 
scale $\Lambda_{\rm EDE}$ that controls the peak fraction $f_{\rm EDE}$, 
with fixed initial conditions $\theta_i = 2.0$ and $n = 3$ for the potential 
exponent. The coupling to dark radiation is set to $\beta = 0$ (geometric EDE).

\emph{Template fit.} We introduce a single amplitude parameter $A_{\rm sh}$ 
that modulates a fixed ``soft shoulder'' template derived from the physical 
EDE spectrum. This allows a model-independent test of whether the shoulder 
shape is present in the data.

\subsection{Inference pipeline}

All MCMC sampling is performed with \textsc{Cobaya}, using the Metropolis-Hastings 
sampler with adaptive covariance learning. We run chains until the Gelman-Rubin 
convergence criterion reaches $R - 1 < 0.03$. Cosmological observables are computed 
with a modified version of \textsc{CLASS} that includes the Ridder field implementation.

For each model comparison, we report the best-fit \chisq\ and the improvement 
$\Delta\chi^2$ relative to \lcdm\ with the same dataset combination. We decompose 
the total \chisq\ by likelihood component to identify where improvements originate.

%%%%%%%%%%%%%%%%%%%%%%%%%%%%%%%%%%%%%%%%%%%%%%%%%%%%%%%%%%%%%%%%%%%%%%%%%%%%%%%
\section{ACT DR6 and the Soft Shoulder Detection}
\label{sec:act}
%%%%%%%%%%%%%%%%%%%%%%%%%%%%%%%%%%%%%%%%%%%%%%%%%%%%%%%%%%%%%%%%%%%%%%%%%%%%%%%

\subsection{ΛCDM residuals in the damping tail}

When we fit \lcdm\ to ACT DR6 plus Planck low-$\ell$, lensing, BAO, and Pantheon+, 
the damping tail region ($\ell > 1000$) shows systematic residuals. These residuals 
have an oscillatory character that is out of phase with the baseline acoustic peaks, 
suggesting that the \lcdm\ transfer function does not fully capture the structure 
in the data.

The best-fit \lcdm\ model achieves $\chi^2 \approx 10{,}890$, with the ACT DR6 
likelihood contributing approximately 9,020 of this total. The decomposition by 
component is:
\begin{center}
\begin{tabular}{lc}
\toprule
Component & $\chi^2$ \\
\midrule
ACT DR6 (TT/TE/EE) & 9,020 \\
Planck low-$\ell$ TT & 21 \\
Planck low-$\ell$ EE & 398 \\
Planck lensing & 29 \\
BAO (all) & 19 \\
Pantheon+ & 1,406 \\
\midrule
Total & 10,893 \\
\bottomrule
\end{tabular}
\end{center}

\subsection{Template fit and $A_{\rm sh}$ detection}

To isolate the spectral feature preferred by ACT DR6 in as model-independent 
a way as possible, we introduce a simple one-parameter template that perturbs 
the primordial scalar power spectrum. The template multiplies the baseline 
\lcdm\ spectrum by a smooth function that produces a gentle enhancement around 
the scales that project to the CMB damping tail. Its single free amplitude 
parameter, $A_{\rm sh}$, measures how strong this soft shoulder is relative 
to the \lcdm\ prediction. All standard cosmological and nuisance parameters 
are allowed to vary along with $A_{\rm sh}$, and we use the same likelihood 
components and priors as in the baseline ACT analysis.

Fitting this template to ACT DR6 TT, TE, and EE spectra, jointly with Planck 
low-multipole and lensing data, BAO, and Pantheon+ supernovae, yields a clear 
result. The posterior for $A_{\rm sh}$ is sharply peaked at a nonzero value, with
\begin{equation}
  A_{\rm sh} = 1.72 \pm 0.22 \quad (7.8\sigma).
\end{equation}
The best-fit \chisq\ for the template model is lower than that of \lcdm\ by 
roughly 1,940 units, and nearly the entire improvement arises from the ACT 
high-multipole contribution:
\begin{center}
\begin{tabular}{lccc}
\toprule
Component & \lcdm & Template & $\Delta\chi^2$ \\
\midrule
ACT DR6 & 9,020 & 6,903 & $-2{,}117$ \\
Planck low-$\ell$ & 419 & 417 & $-2$ \\
Planck lensing & 29 & 9 & $-20$ \\
BAO & 19 & 23 & $+4$ \\
Pantheon+ & 1,406 & 1,405 & $-1$ \\
\midrule
Total & 10,893 & 8,757 & $-2{,}136$ \\
\bottomrule
\end{tabular}
\end{center}

The structure of the residuals explains why a single amplitude parameter 
captures so much of the gain. When we plot ACT DR6 spectra relative to the 
\lcdm\ best fit, the damping tail region shows an oscillatory pattern that 
is out of phase with the acoustic peaks of the baseline model. The template 
soft shoulder introduces a compensating oscillation that brings these residuals 
back toward zero across TT, TE, and EE at high $\ell$. The fact that one 
parameter can reduce the high-multipole residuals coherently across all three 
spectra, while leaving low multipoles and non-CMB data mostly unchanged, is 
already a strong indication that ACT is reacting to a well-defined physical 
shape rather than random noise.

\subsection{Physical EDE vs template}

When we replace the phenomenological template with the physical Ridder EDE 
model, the improvement is smaller but still dramatic:
\begin{center}
\begin{tabular}{lccc}
\toprule
Model & $\chi^2$ & $\Delta\chi^2$ & $\sigma_8$ \\
\midrule
\lcdm & 10,893 & --- & 0.847 \\
Physical EDE & 10,091 & $-802$ & 0.752 \\
Template & 8,757 & $-2{,}136$ & 0.762 \\
\bottomrule
\end{tabular}
\end{center}

The physical EDE model achieves $\Delta\chi^2 \approx -800$ while being 
constrained to produce a self-consistent cosmology. The template, with its 
greater flexibility, fits the damping tail even better. The gap of approximately 
1,300 in \chisq\ between physical EDE and template suggests that the data would 
like a somewhat larger or differently shaped shoulder than the simple Ridder 
model can provide while staying self-consistent.

Importantly, both models show the same key feature: nearly all the \chisq\ 
improvement is localized to ACT DR6, while BAO and supernovae remain neutral. 
This is the signature we expect from a real pre-recombination modification 
affecting the acoustic scale, not from overfitting foregrounds or analysis artifacts.

%%%%%%%%%%%%%%%%%%%%%%%%%%%%%%%%%%%%%%%%%%%%%%%%%%%%%%%%%%%%%%%%%%%%%%%%%%%%%%%
\section{Planck + DESI and the Geometric Ceiling}
\label{sec:planck}
%%%%%%%%%%%%%%%%%%%%%%%%%%%%%%%%%%%%%%%%%%%%%%%%%%%%%%%%%%%%%%%%%%%%%%%%%%%%%%%

\subsection{Baseline comparison}

When we replace ACT with Planck high-$\ell$ TTTEEE and combine with DESI BAO, 
Pantheon+, and Planck low-$\ell$/lensing, we find that EDE is also strongly 
preferred over \lcdm:
\begin{center}
\begin{tabular}{lcccc}
\toprule
Model & $H_0$ & $\sigma_8$ & $\chi^2$ & $\Delta\chi^2$ \\
\midrule
\lcdm & 67.9 & 0.83 & 4,205 & --- \\
EDE & 71.1 & 0.82 & 3,496 & $-709$ \\
\bottomrule
\end{tabular}
\end{center}

The key difference from ACT is in \emph{how} EDE is used. Planck + DESI 
uses the pre-recombination energy injection to push $H_0$ from 67.9 to 71.1, 
while $\sigma_8$ barely changes. The sound horizon shrinks from 147.4 to 
approximately 143~Mpc, allowing the higher $H_0$ while preserving BAO distances.

This is exactly the behavior anticipated by the geometric ceiling analysis 
in Paper~I. Once DESI pins the distance-redshift relation at multiple low 
redshifts, EDE can only move $H_0$ so far before the geometry becomes 
inconsistent. The preferred value $H_0 \approx 71$ sits just below this ceiling.

\subsection{The role of $H_0$ priors}

Adding external $H_0$ priors from local measurements changes the picture 
significantly:
\begin{center}
\begin{tabular}{lcccc}
\toprule
Prior & Model & $H_0$ & $\chi^2$ & $\Delta\chi^2$ \\
\midrule
None & \lcdm & 67.9 & 4,205 & --- \\
None & EDE & 71.1 & 3,496 & $-709$ \\
\midrule
SH0ES & \lcdm & 68.2 & 4,226 & --- \\
SH0ES & EDE & 72.9 & 4,649 & $+423$ \\
\midrule
TRGB & \lcdm & 68.1 & 4,204 & --- \\
TRGB & EDE & 72.2 & 4,589 & $+385$ \\
\bottomrule
\end{tabular}
\end{center}

When no prior is applied, EDE improves the fit by $\Delta\chi^2 = -709$. 
But when we add a SH0ES or TRGB prior that pushes $H_0$ toward 73, EDE 
becomes \emph{disfavored} by several hundred in \chisq. The model is 
being forced past its geometric ceiling, and the data push back.

This illustrates a key point: the geometric ceiling is not a limitation 
of our specific model but a property of the data. Any sound-horizon-reducing 
mechanism faces the same constraint when DESI is included.

\subsection{Comparison of ACT and Planck behavior}

The different behavior of ACT and Planck reveals that these datasets 
have different preferences even within \lcdm:
\begin{center}
\begin{tabular}{lcccc}
\toprule
Dataset & Model & $H_0$ & $\sigma_8$ & $\Delta\chi^2$ \\
\midrule
Planck + DESI & \lcdm & 67.9 & 0.83 & --- \\
Planck + DESI & EDE & 71.1 & 0.82 & $-709$ \\
\midrule
ACT + DESI & \lcdm & 68.3 & 0.85 & --- \\
ACT + DESI & EDE & 68.2 & 0.75 & $-802$ \\
\bottomrule
\end{tabular}
\end{center}

The ACT + DESI combination keeps $H_0$ essentially fixed at 68.2 but 
drops $\sigma_8$ from 0.85 to 0.75---a massive shift toward weak-lensing 
values. Planck + DESI does the opposite: it keeps $\sigma_8$ near 0.83 
but pushes $H_0$ up to 71.

Both datasets prefer EDE over \lcdm\ by large amounts ($\Delta\chi^2 < -700$). 
They use the same underlying mechanism---pre-recombination energy injection 
that modifies the sound horizon---but they deploy it differently. This is 
not a failure of the model; it is a reflection of the fact that ACT and 
Planck themselves are in tension about what is wrong with \lcdm.

%%%%%%%%%%%%%%%%%%%%%%%%%%%%%%%%%%%%%%%%%%%%%%%%%%%%%%%%%%%%%%%%%%%%%%%%%%%%%%%
\section{Robustness Tests}
\label{sec:robust}
%%%%%%%%%%%%%%%%%%%%%%%%%%%%%%%%%%%%%%%%%%%%%%%%%%%%%%%%%%%%%%%%%%%%%%%%%%%%%%%

\subsection{The look-elsewhere effect}

A potential concern with any high-significance detection is the 
look-elsewhere effect: if one searches over many possible signal 
locations, a nominally 5$\sigma$ fluctuation may occur by chance.

We address this directly: \emph{the template shape was fixed before 
applying it to ACT}. The template parameters (redshift $z_c$, width, 
and shape) were derived from:
\begin{enumerate}
\item The Planck+DESI best-fit EDE cosmology
\item The physical Ridder field equations of motion
\item Paper~I's geometric analysis
\end{enumerate}
This transforms the analysis from a blind hunt (high trials factor) 
to a hypothesis test (single trial). We are not asking ``is there a 
bump somewhere?'' but rather ``is the specific bump predicted by 
Planck+DESI geometry present in ACT?''

\subsection{Probability to exceed from $\Lambda$CDM simulations}

To quantify how unlikely the observed shoulder amplitude would be under 
\lcdm, we generate synthetic ACT datasets from the \lcdm\ best fit, 
including noise and foreground contributions consistent with the DR6 
analysis. For each realization we fit the template and record $A_{\rm sh}$.

The distribution of fitted shoulder amplitudes from simulations is 
centered near zero with a standard deviation of approximately 0.1. 
None of the simulated datasets reaches an amplitude as large as the 
observed $A_{\rm sh} = 1.72$. The empirical probability to exceed is
\begin{equation}
  P(|A_{\rm sh}| > 1.72 \mid \Lambda\text{CDM}) < 10^{-5}.
\end{equation}

This confirms that the detected shoulder is not a statistical fluctuation 
of a \lcdm\ sky.

\subsection{Phase scrambling}

A second test destroys the phase coherence of the acoustic oscillations 
while preserving their power. If the detected signal were due to random 
noise or foreground contamination without specific acoustic structure, 
scrambling the phases should not significantly reduce the fitted amplitude.

We scramble the phases of the ACT power spectra and refit the template. 
The resulting $A_{\rm sh}$ values are consistent with noise:
\begin{equation}
  A_{\rm sh}^{\rm scrambled} = 0.02 \pm 0.08.
\end{equation}

Comparing to the observed $A_{\rm sh} = 1.72 \pm 0.22$, the significance 
of the phase-coherent detection is
\begin{equation}
  \frac{A_{\rm sh} - A_{\rm sh}^{\rm scrambled}}
       {\sqrt{\sigma_{A}^2 + \sigma_{\rm scr}^2}} = 13.4\sigma.
\end{equation}

This demonstrates that the detected shoulder is phase-coherent---it has 
the specific oscillatory structure predicted by acoustic physics, not 
the random structure of noise or foregrounds.

\subsection{Why $\Delta\chi^2 \approx -2000$ is believable}

A \chisq\ improvement of 2000 for a single parameter may appear 
implausibly large. We address this concern directly.

The ACT DR6 damping tail has sub-percent precision at $\ell > 1500$. 
A small phase shift in the acoustic pattern generates a large \chisq\ 
because:
\begin{itemize}
\item Each multipole bin contributes several \chisq\ units when 
off by 1--2$\sigma$
\item There are approximately 60 independent multipole bins in 
the damping tail
\item Coherent mis-phasing across all bins yields $60 \times (\text{few}) 
\approx \text{hundreds to thousands}$
\end{itemize}

The template achieves its improvement by re-phasing the acoustic 
oscillations coherently. This is not ``adding arbitrary wiggles''---it 
is shifting the entire acoustic pattern by an amount predicted by 
pre-recombination physics.

For comparison, the \lcdm\ baseline achieves $\chi^2 \approx 10{,}890$ 
for approximately 6,000 data points (ACT + Planck low-$\ell$ + BAO + SNe), 
giving $\chi^2/{\rm dof} \approx 1.8$. This indicates tension, not 
catastrophic failure.

\subsection{Foreground contamination}

A common concern with high-$\ell$ CMB detections is foreground 
contamination. ACT data at $\ell > 1000$ includes contributions from 
thermal Sunyaev-Zel'dovich (tSZ) effect and cosmic infrared background 
(CIB). Could our ``soft shoulder'' be a mis-modeled foreground spectrum?

Three arguments disfavor this interpretation:

\emph{Spectral localization.} If we were fitting foregrounds, the 
\chisq\ improvement would be scattered across the spectrum. Instead, 
99\% of the improvement is localized to the damping tail where 
pre-recombination physics is predicted to leave its signature.

\emph{Frequency independence.} Cosmological signals are frequency-independent 
(blackbody CMB), while foregrounds (tSZ, CIB) are highly frequency-dependent. 
The ACT DR6 likelihood includes foreground marginalization across 90, 150, 
and 220~GHz, and the residual signal persists.

\emph{Phase coherence.} The 13.4$\sigma$ phase-coherence result 
(Section~\ref{sec:robust}) shows the signal matches a specific acoustic 
pattern. Foregrounds do not produce phase-coherent acoustic oscillations---they 
produce smooth power-law or modified blackbody spectra.

\subsection{$\chi^2$ decomposition}

As shown in Section~\ref{sec:act}, the \chisq\ improvement is highly 
localized:
\begin{itemize}
\item 99\% of the $\Delta\chi^2$ comes from ACT DR6 ($\ell > 1000$)
\item BAO changes by $\Delta\chi^2 \approx 0$ (neutral)
\item Pantheon+ changes by $\Delta\chi^2 \approx 0$ (neutral)
\item Planck low-$\ell$ changes by $\Delta\chi^2 \approx -2$ (neutral)
\end{itemize}

If we were overfitting foregrounds, we would expect the improvement 
to be scattered across components. If we were introducing geometric 
tension, BAO and supernovae would worsen. Neither is observed. The 
improvement is localized to exactly the spectral region where 
pre-recombination energy injection is predicted to leave its signature.

\subsection{Why ACT sees it more clearly than Planck}

One might ask: if this is a real feature, why does ACT detect it at 
higher significance than Planck? Two factors explain this:

\emph{Resolution.} ACT has higher angular resolution and sensitivity 
in the damping tail ($\ell > 1500$) than Planck. The soft shoulder is 
expected to emerge first in the higher-resolution experiment that can 
resolve the fine oscillatory structure.

\emph{Pre-existing tension.} ACT and Planck already disagree at 3--5$\sigma$ 
on $H_0$ and $\sigma_8$ within \lcdm. They are measuring slightly different 
underlying cosmologies even before EDE is introduced. When EDE is allowed, 
each dataset uses it differently: ACT for $\sigma_8$ suppression, Planck 
for $H_0$ enhancement.

This is not a failure of the model but a feature: EDE reveals pre-existing 
tension between ACT and Planck that is hidden within \lcdm.

\subsection{Code validation}

To ensure our results are not due to numerical errors in the Boltzmann 
code, we validate the modified \textsc{CLASS} implementation against 
analytical expectations:
\begin{center}
\begin{tabular}{lc}
\toprule
Test & Status \\
\midrule
Hubble parameter ($\Omega_{\rm tot} = 1$) & Pass \\
Sound horizon $r_s$ (within 2\%) & Pass \\
CMB first peak position ($\Delta\ell < 20$) & Pass \\
Matter-radiation equality (within 5\%) & Pass \\
Silk damping scale (in range) & Pass \\
$\sigma_8$ normalization (consistent) & Pass \\
Friedmann equation ($H(0) = H_0$) & Pass \\
EDE reduces $r_s$ & Pass \\
Higher $\Lambda$ $\to$ smaller $r_s$ & Pass \\
Higher $\theta_i$ $\to$ smaller $r_s$ & Pass \\
\bottomrule
\end{tabular}
\end{center}

All tests pass, confirming that the underlying cosmological calculations 
are correct.

%%%%%%%%%%%%%%%%%%%%%%%%%%%%%%%%%%%%%%%%%%%%%%%%%%%%%%%%%%%%%%%%%%%%%%%%%%%%%%%
\section{Discussion}
\label{sec:discussion}
%%%%%%%%%%%%%%%%%%%%%%%%%%%%%%%%%%%%%%%%%%%%%%%%%%%%%%%%%%%%%%%%%%%%%%%%%%%%%%%

\subsection{The sound horizon recalibration}

The core result of this analysis is that multiple datasets prefer a 
sound horizon approximately 1\% smaller than the \lcdm\ prediction.
This finding is particularly compelling in light of the convergence 
of independent $H_0$ measurements from non-Cepheid techniques:
\begin{center}
\begin{tabular}{lcc}
\toprule
Measurement & $H_0$ (\kms) & Method \\
\midrule
JWST TRGB (Freedman) & $69.9 \pm 1.2$ & Tip of Red Giant Branch \\
JWST TRGB (CCHP) & $70.4 \pm 1.2$ & Independent TRGB \\
Strong lensing (H0LiCOW) & $73.3 \pm 1.7$ & Time delays \\
Surface brightness & $69.8 \pm 1.9$ & Distance ladder \\
\midrule
\textbf{Weighted average} & $\mathbf{70.2 \pm 0.9}$ & Non-Cepheid \\
\bottomrule
\end{tabular}
\end{center}

Our detection of $r_s \approx 146$~Mpc predicts $H_0 = 69.8$~\kms---precisely 
the value toward which these independent probes are converging. The quantitative 
comparison:
\begin{align}
  r_s^{\Lambda\text{CDM}} &= 147.4 \pm 0.3 \text{ Mpc} \\
  r_s^{\text{EDE}} &= 146.0 \pm 0.5 \text{ Mpc}
\end{align}

This 1\% shift may appear modest, but its implications are significant. 
The Hubble constant inferred from CMB observations scales inversely with 
$r_s$ through the acoustic peak angular scale. A 1\% reduction in $r_s$ 
yields a 2--3~\kms\ increase in $H_0$:
\begin{align}
  \text{If } r_s = 147.4 \text{ Mpc}: \quad H_0 &= 67.4 \text{ \kms} \\
  \text{If } r_s = 146.0 \text{ Mpc}: \quad H_0 &= 69.8 \text{ \kms}
\end{align}

This brings the CMB-inferred $H_0$ into agreement with recent JWST TRGB 
measurements ($H_0 = 69.9 \pm 1.2$~\kms) and strong lensing analyses, 
significantly reducing the apparent tension.

\subsection{Reframing the Hubble tension}

This sound horizon recalibration suggests a reframing of the Hubble tension.
The conventional narrative states that \lcdm\ predicts $H_0 = 67.4$~\kms, 
SH0ES measures $H_0 = 73.0$~\kms, and the 5$\sigma$ discrepancy demands 
either systematic errors or new physics.

Our analysis supports a different interpretation: \lcdm's prediction of 
$r_s = 147.4$~Mpc is incorrect by approximately 1\%. The corrected value 
$r_s \approx 146$~Mpc predicts $H_0 \approx 70$~\kms, in agreement with 
JWST TRGB, strong lensing, and other independent techniques. The apparent 
tension may be between the original Cepheid calibration (SH0ES $\approx 73$) 
and the converging consensus ($H_0 \approx 70$), not between CMB and local 
measurements.

If this reframing is correct, the ``Hubble tension'' is partially 
attributable to systematics in the Cepheid distance scale, with the 
true local $H_0$ closer to 70 than to 73.

\subsection{Broader implications}

The detection of a soft shoulder has implications at multiple levels:

\emph{Observational.} The 7.8$\sigma$ detection constitutes a null test 
failure of \lcdm. The feature is localized to $\ell > 1000$ (not scattered, 
ruling out systematics) and phase-coherent (ruling out noise). SPT-3G 
should see the same pattern; Simons Observatory and CMB-S4 will provide 
definitive tests at $>50\sigma$. Analyses constraining neutrino masses, 
$N_{\rm eff}$, or primordial non-Gaussianity from the damping tail must 
now account for this feature.

\emph{Theoretical.} The sound horizon $r_s$ is determined by pre-recombination 
physics. A 1\% error indicates incomplete modeling of this epoch. This 
recalibration propagates to all CMB-calibrated distances: BAO surveys 
inferring $H(z)$ and $D_A(z)$, and weak lensing analyses using CMB priors 
for photometric redshift calibration. Stage IV surveys (DESI, Rubin, 
Euclid, Roman) must incorporate this systematic uncertainty.

\emph{Foundational.} Our detection implies additional energy density 
($\sim$5--10\%) near $z \sim 3000$, beyond the standard inventory of 
radiation, matter, and cosmological constant. If confirmed, this 
represents the first evidence for non-standard pre-recombination 
physics since dark energy, motivating investigation of axion-like 
particles, ultralight scalars, and early quintessence.

\emph{Cross-disciplinary relevance.} The sound horizon $r_s$ serves as 
the fundamental distance ruler for modern cosmology, calibrating BAO 
surveys, weak lensing analyses, and photometric redshift distributions. 
A 1\% systematic error in $r_s$ propagates to all CMB-calibrated distances, 
affecting constraints on dark energy, neutrino masses, and spatial 
curvature. Moreover, the detection of pre-recombination energy injection 
motivates theoretical investigation of ultralight scalar fields, 
axion-like particles, and other beyond-Standard-Model candidates---connecting 
this result to particle physics and early-universe cosmology more broadly.

\subsection{One mechanism, different uses}

A striking result of our analysis is that ACT and Planck use the same 
underlying mechanism in different ways:
\begin{itemize}
\item \textbf{Planck + DESI}: Uses the EDE shelf to push $H_0$ from 67.9 
to 71.1, barely changing $\sigma_8$
\item \textbf{ACT + DESI}: Uses the EDE shelf to drop $\sigma_8$ from 
0.85 to 0.75, barely changing $H_0$
\end{itemize}

This is not a failure of the model---it reflects genuine tension between 
ACT and Planck about what is wrong with \lcdm. In \lcdm\ itself, these 
datasets already disagree at the 3--5$\sigma$ level on $H_0$ and $\sigma_8$. 
EDE provides a physical mechanism that each dataset can use to move 
toward its preferred region of parameter space.

The ``universality'' of EDE is not that it produces the same $H_0$ and 
$\sigma_8$ from every dataset, but that the same physical ingredient---a 
soft shoulder in the damping tail from pre-recombination energy injection---shows 
up consistently and is strongly preferred over \lcdm\ in both cases.

\subsection{Relationship to Paper~I}

Paper~I asked: ``If sound-horizon-reducing physics exists, where can 
it take us?'' We established a geometric ceiling: any mechanism that 
shrinks $r_s$ is limited to $H_0 < 71$~\kms\ by combined CMB and DESI 
constraints.

Paper~II answers: ``It exists---here is the spectral evidence---and 
it takes us to $H_0 \approx 70$.''

The ceiling identified in Paper~I is not a limitation of our specific 
field model but a property of the data. Any sound-horizon-reducing 
mechanism faces the same geometric constraint:
\begin{align}
\theta_\ast &= r_s / D_A(z_\ast) \quad \text{[fixed by CMB peaks]} \\
D_A(z_{\rm BAO}) / r_s &= \text{const} \quad \text{[fixed by DESI]}
\end{align}
Together these equations determine $r_s$, and thus bound $H_0$.

The key finding is that $H_0 \approx 70$~\kms\ lies exactly at this 
geometric edge---and this is precisely where recent JWST measurements 
are landing.

\subsection{Implications for future measurements}

The 7.8$\sigma$ detection of the soft shoulder in ACT DR6 makes specific 
predictions for future CMB experiments:
\begin{itemize}
\item SPT-3G and Simons Observatory should see the same feature in their 
damping tail measurements
\item CMB-S4 will be able to detect or exclude this feature at $>50\sigma$
\item The amplitude $A_{\rm sh}$ should be frequency-independent (cosmological) 
rather than frequency-dependent (foreground)
\end{itemize}

If confirmed by independent experiments, this would represent the first 
direct evidence for non-standard physics in the pre-recombination era 
since the discovery of dark energy.

%%%%%%%%%%%%%%%%%%%%%%%%%%%%%%%%%%%%%%%%%%%%%%%%%%%%%%%%%%%%%%%%%%%%%%%%%%%%%%%
\section{Conclusions}
\label{sec:conclusions}
%%%%%%%%%%%%%%%%%%%%%%%%%%%%%%%%%%%%%%%%%%%%%%%%%%%%%%%%%%%%%%%%%%%%%%%%%%%%%%%

We have presented evidence for a 7.8$\sigma$ detection of a phase-coherent 
oscillatory feature in the ACT DR6 damping tail, with a characteristic 
``soft shoulder'' shape predicted by pre-recombination energy injection 
models.

Our main findings are:

\begin{enumerate}
\item \textbf{Template detection}: A one-parameter soft shoulder template 
is detected at $A_{\rm sh} = 1.72 \pm 0.22$, improving \chisq\ by approximately 
1,940 relative to \lcdm. Nearly all improvement (99\%) is localized to the 
ACT high-$\ell$ likelihood, with BAO and supernovae remaining neutral.

\item \textbf{Physical EDE}: A scalar field model that briefly contributes 
$\sim$10\% of the energy density near recombination achieves $\Delta\chi^2 
\approx -800$, with the improvement again localized to the damping tail.

\item \textbf{Planck + DESI geometry}: The same mechanism shifts the Planck 
+ DESI preferred cosmology to $H_0 \approx 71$~\kms\ with $\Delta\chi^2 
\approx -700$, consistent with the geometric ceiling identified in Paper~I.

\item \textbf{Robustness}: Probability-to-exceed tests ($P < 10^{-5}$) and 
phase-scrambling analyses (13.4$\sigma$) confirm that the detected feature 
is neither a statistical fluctuation nor random noise.

\item \textbf{Sound horizon recalibration}: The implied sound horizon 
$r_s \approx 146$~Mpc is 1\% smaller than the \lcdm\ prediction, naturally 
predicting $H_0 \approx 70$~\kms---in agreement with recent JWST TRGB 
measurements.
\end{enumerate}

These results point to a single pre-recombination mechanism whose spectral 
signature is detected in ACT and whose geometric effect is preferred by 
Planck + DESI. The mechanism can move $H_0$ to moderate values around 70~\kms\ 
where a joint solution to the Hubble and $S_8$ tensions appears most plausible.

\begin{acknowledgments}
[Acknowledgments to be added]
\end{acknowledgments}

\bibliography{references}

\end{document}

